% Options for packages loaded elsewhere
\PassOptionsToPackage{unicode}{hyperref}
\PassOptionsToPackage{hyphens}{url}
\documentclass[
  ignorenonframetext,
  aspectratio=169,
]{beamer}
\newif\ifbibliography
\usepackage{pgfpages}
\setbeamertemplate{caption}[numbered]
\setbeamertemplate{caption label separator}{: }
\setbeamercolor{caption name}{fg=normal text.fg}
\beamertemplatenavigationsymbolsempty
% remove section numbering
\setbeamertemplate{part page}{
  \centering
  \begin{beamercolorbox}[sep=16pt,center]{part title}
    \usebeamerfont{part title}\insertpart\par
  \end{beamercolorbox}
}
\setbeamertemplate{section page}{
  \centering
  \begin{beamercolorbox}[sep=12pt,center]{section title}
    \usebeamerfont{section title}\insertsection\par
  \end{beamercolorbox}
}
\setbeamertemplate{subsection page}{
  \centering
  \begin{beamercolorbox}[sep=8pt,center]{subsection title}
    \usebeamerfont{subsection title}\insertsubsection\par
  \end{beamercolorbox}
}
% Prevent slide breaks in the middle of a paragraph
\widowpenalties 1 10000
\raggedbottom
\AtBeginPart{
  \frame{\partpage}
}
\AtBeginSection{
  \ifbibliography
  \else
    \frame{\sectionpage}
  \fi
}
\AtBeginSubsection{
  \frame{\subsectionpage}
}
\usepackage{iftex}
\ifPDFTeX
  \usepackage[T1]{fontenc}
  \usepackage[utf8]{inputenc}
  \usepackage{textcomp} % provide euro and other symbols
\else % if luatex or xetex
  \usepackage{unicode-math} % this also loads fontspec
  \defaultfontfeatures{Scale=MatchLowercase}
  \defaultfontfeatures[\rmfamily]{Ligatures=TeX,Scale=1}
\fi
\usepackage{lmodern}
\ifPDFTeX\else
  % xetex/luatex font selection
\fi
% Use upquote if available, for straight quotes in verbatim environments
\IfFileExists{upquote.sty}{\usepackage{upquote}}{}
\IfFileExists{microtype.sty}{% use microtype if available
  \usepackage[]{microtype}
  \UseMicrotypeSet[protrusion]{basicmath} % disable protrusion for tt fonts
}{}
\makeatletter
\@ifundefined{KOMAClassName}{% if non-KOMA class
  \IfFileExists{parskip.sty}{%
    \usepackage{parskip}
  }{% else
    \setlength{\parindent}{0pt}
    \setlength{\parskip}{6pt plus 2pt minus 1pt}}
}{% if KOMA class
  \KOMAoptions{parskip=half}}
\makeatother
\usepackage{longtable,booktabs,array}
\usepackage{caption}
\captionsetup[table]{skip=6pt}
\newcounter{none} % for unnumbered tables
\usepackage{calc} % for calculating minipage widths
\usepackage{caption}
% Make caption package work with longtable
\makeatletter
\def\fnum@table{\tablename~\thetable}
\makeatother
\setlength{\emergencystretch}{3em} % prevent overfull lines
\providecommand{\tightlist}{%
  \setlength{\itemsep}{0pt}\setlength{\parskip}{0pt}}
% Auto-generated by infrastructure/rendering/_pdf_latex_helpers.py::extract_math_font_preamble.
% Minimal header passed to Pandoc via -H so Beamer slides pick up the same math font
% as the combined-PDF path. Adjust the manuscript's preamble.md to change the font globally.
\usepackage{unicode-math}
\setmathfont{latinmodern-math.otf}

% Manuscript macro/package fallbacks (newcommand -> providecommand).
% Core mathematics
\usepackage{amsmath}
\usepackage{amssymb}
% Document layout
% Tables
\usepackage{booktabs}
% Caption styling (booktabs already loaded above). Small caption font with a
% bold label ("Figure 1", "Table 2") gives the math/figure-heavy layout a
% consistent, journal-like caption rhythm without fighting pandoc-crossref —
% pandoc emits standard \caption{}, and caption/captionsetup only restyle it.
% ── Theorem environments ─────────────────────────────────────────────────
% amsthm is loaded above. The FedGVI<->Friston bridge is stated as a small
% number of theorems/definitions/lemmas/propositions/corollaries; sharing one
% counter (definition/lemma/proposition/corollary all step the `theorem`
% counter) gives a single monotone numbering 1, 2, 3, ... across all kinds, so
% authors can reference them as Theorem 1, Definition 2, Corollary 3 without a
% per-kind counter drifting out of order. Plain style for theorem-like results,
% definition style (upright body) for definitions.
% (algorithm2e intentionally not loaded: this manuscript states results as
% numbered amsthm Theorems/Definitions, not algorithm2e pseudocode.)
% Code listings
\usepackage{listings}
% Typography and formatting
% (siunitx intentionally not loaded: no \SI/\num units appear in this manuscript.)
% Cross-references and citations
% (cleveref and titlesec intentionally not loaded: unavailable in this TeX tree
% and not required. Cross-references resolve through pandoc-crossref's own
% \ref-style names — no \cref appears — and section heading styling falls back
% to the pandoc/LaTeX defaults.)
% ── TeX Live Basic-compatible mono font for code listings ────────────
% Latin Modern Mono is bundled with TeX Live Basic and keeps the manuscript
% renderable on minimal TeX installations. Mathematical Unicode coverage is
% provided separately by Latin Modern Math below, so the code font does not
% need a private JuliaMono installation.
%
% Use the TeX font filename rather than the system-family name: the minimal
% TeX Live fontconfig database may not register the OpenType family even though
% kpsewhich can resolve the bundled files.
% Math font for unicode-math: Latin Modern Math (TeX Live) has full BMP
% coverage including U+2223 (\mid), U+226A/226B (\ll/\gg), and the Greek/
% blackboard letters used in equations. Without an explicit \setmathfont,
% unicode-math falls back to lmroman text font which lacks several glyphs
% and emits "Missing character" warnings on every \mid in math mode.

% Slide layout defaults for warning-clean scientific decks.
\usepackage{etoolbox}
\IfFileExists{xurl.sty}{\usepackage{xurl}}{}
\IfFileExists{seqsplit.sty}{\usepackage{seqsplit}}{\newcommand{\seqsplit}[1]{#1}}
\protected\def\breakseq#1{\seqsplit{#1}}
\protected\def\breaktt#1{\begingroup\ttfamily\seqsplit{#1}\endgroup}
\setlength{\emergencystretch}{6em}
\tolerance=5000
\hbadness=10000
\hfuzz=1pt
\setlength{\tabcolsep}{2pt}
\AtBeginEnvironment{longtable}{\tiny\renewcommand{\arraystretch}{0.86}\setlength{\tabcolsep}{1pt}}
\AtBeginEnvironment{tabular}{\tiny\renewcommand{\arraystretch}{0.86}\setlength{\tabcolsep}{1pt}}
\AtBeginEnvironment{equation}{\tiny}
\AtBeginEnvironment{equation*}{\tiny}
\AtBeginEnvironment{align}{\tiny}
\AtBeginEnvironment{align*}{\tiny}
\AtBeginEnvironment{itemize}{\footnotesize}
\AtBeginEnvironment{enumerate}{\footnotesize}
\AtBeginEnvironment{description}{\footnotesize}
\setbeamerfont{caption}{size=\tiny}
\setbeamerfont{caption name}{size=\tiny}
\setbeamerfont{normal text}{size=\small}
\setbeamerfont{frametitle}{size=\small}
\setbeamerfont{section title}{size=\footnotesize}
\setbeamerfont{subsection title}{size=\footnotesize}
\setbeamertemplate{section page}{%
  \centering
  \begin{beamercolorbox}[sep=12pt,center,wd=\paperwidth]{section title}
    \parbox{0.86\paperwidth}{\centering\usebeamerfont{section title}\insertsection\par}
  \end{beamercolorbox}
}
\setbeamertemplate{subsection page}{%
  \centering
  \begin{beamercolorbox}[sep=8pt,center,wd=\paperwidth]{subsection title}
    \parbox{0.86\paperwidth}{\centering\usebeamerfont{subsection title}\insertsubsection\par}
  \end{beamercolorbox}
}
\setlength{\abovecaptionskip}{2pt}
\setlength{\belowcaptionskip}{0pt}

% Natbib and cross-reference fallbacks — slides don't load natbib
% or cleveref, but combined-PDF manuscript prose may emit these
% commands. The fallback renders citations as a bracketed key list
% and cross-references as detokenized labels so slides don't fail on
% undefined control sequences or raw underscores. \providecommand is
% a no-op if packages load later.
\providecommand{\citep}[1]{[#1]}
\providecommand{\citet}[1]{#1}
\providecommand{\citealp}[1]{#1}
\providecommand{\citeauthor}[1]{#1}
\providecommand{\citeyear}[1]{#1}
\providecommand{\cref}[1]{\texttt{\detokenize{#1}}}
\providecommand{\Cref}[1]{\texttt{\detokenize{#1}}}
\providecommand{\crefrange}[2]{\texttt{\detokenize{#1}}--\texttt{\detokenize{#2}}}
\providecommand{\Crefrange}[2]{\texttt{\detokenize{#1}}--\texttt{\detokenize{#2}}}
\providecommand{\paragraph}[1]{\textbf{#1}\ }

% Opt-in accessible presentation profile.
\setbeamerfont{normal text}{size*={20pt}{24pt}}
\setbeamerfont{frametitle}{size*={28pt}{32pt}}
\setbeamerfont{section title}{size*={28pt}{32pt}}
\setbeamerfont{subsection title}{size*={28pt}{32pt}}
\setbeamerfont{caption}{size*={16pt}{19pt}}
\setbeamerfont{caption name}{size*={16pt}{19pt}}
\setbeamerfont{itemize/enumerate body}{size*={20pt}{24pt}}
\setbeamerfont{itemize/enumerate subbody}{size*={20pt}{24pt}}
\setbeamerfont{itemize/enumerate subsubbody}{size*={20pt}{24pt}}
\setbeamerfont{itemize item}{size*={20pt}{24pt}}
\setbeamerfont{itemize subitem}{size*={20pt}{24pt}}
\setbeamerfont{itemize subsubitem}{size*={20pt}{24pt}}
\setbeamerfont{description body}{size*={20pt}{24pt}}
\setbeamerfont{description item}{size*={20pt}{24pt}}
\setbeamerfont{quote}{size*={20pt}{24pt}}
\setbeamerfont{quotation}{size*={20pt}{24pt}}
\AtBeginEnvironment{longtable}{\fontsize{20pt}{24pt}\selectfont}
\AtBeginEnvironment{tabular}{\fontsize{20pt}{24pt}\selectfont}
\AtBeginEnvironment{equation}{\fontsize{20pt}{24pt}\selectfont}
\AtBeginEnvironment{equation*}{\fontsize{20pt}{24pt}\selectfont}
\AtBeginEnvironment{align}{\fontsize{20pt}{24pt}\selectfont}
\AtBeginEnvironment{align*}{\fontsize{20pt}{24pt}\selectfont}
\AtBeginEnvironment{itemize}{\fontsize{20pt}{24pt}\selectfont}
\AtBeginEnvironment{enumerate}{\fontsize{20pt}{24pt}\selectfont}
\AtBeginEnvironment{description}{\fontsize{20pt}{24pt}\selectfont}
\AtBeginDocument{\usebeamerfont{normal text}}
\newlength{\AFSlideTableWidth}
% The fixed footer reservation is calibrated with the semantic seven-line regular-body budget.
\setbeamertemplate{footline}{%
  \leavevmode\vbox{%
    \hbox{%
      \begin{beamercolorbox}[wd=\paperwidth,ht=16pt,dp=3pt,center]{author in head/foot}%
        \fontsize{16pt}{19pt}\selectfont Untagged PDF derivative \textbar\ \href{../web/index.html}{HTML reader}%
      \end{beamercolorbox}%
    }%
    \vskip6pt%
  }%
}
% Accessible footer reserved height: 25pt.

\newtheorem{proposition}{Proposition}
\newtheorem{hypothesis}{Hypothesis}
\newtheorem{remark}[theorem]{Remark}
\newtheorem{axiom}{Axiom}
\newtheorem{property}{Property}
\usepackage{bookmark}
\IfFileExists{xurl.sty}{\usepackage{xurl}}{} % add URL line breaks if available
\urlstyle{same}
% fallback for those not using the hyperref driver hyperxmp:
\makeatletter
\@ifundefined{xmpquote}{\newcommand{\xmpquote}[1]{#1}}{}
\makeatother
\hypersetup{
  hidelinks,
  pdfcreator={LaTeX via pandoc}}

\author{\texorpdfstring{}{}}
\date{}

\begin{document}

\begin{frame}{Future work: testing the open boundaries}
\protect\phantomsection\label{sec:future}
The staged research program follows directly from the boundaries in
sec.~8.13. It separates public-library reproducibility,
server theory, portable source-protocol work, task generalization,
\end{frame}

\begin{frame}{Future work: testing the\ldots{} (part 2)}
and deployment validation so that success in one lane cannot silently
certify another.
\end{frame}

\begin{frame}{Future work: testing the\ldots{} (part 3)}
The order is deliberate. Simulation-study guidance recommends declaring
the estimand, data-generating mechanism, and Monte Carlo precision
before treating replication as evidence (Morris et al. 2019),
\end{frame}

\begin{frame}{Future work: testing the\ldots{} (part 4)}
while Monte Carlo error should be reported separately from an interval
or a hypothesis test (Koehler et al. 2009).
\end{frame}

\begin{frame}{Future work: testing the\ldots{} (part 5)}
For nested agents, trials, and seeds, the resampling unit must respect
the dependence structure (Loy and Korobova 2021). Accordingly, the phase
plan records a primary unit and a falsifier for each extension; a larger
sample or more elaborate diagram is not itself a stronger claim.
\end{frame}

\begin{frame}{Future work: testing the\ldots{} (part 6)}
The implementation registry also separates smoke, pilot, and
confirmatory profiles. Pilot worlds select budgets and calibration
settings but never enter confirmatory intervals or headline values.
\end{frame}

\begin{frame}{Future work: testing the\ldots{} (part 7)}
Each completed run must bind its source bundle, configuration, dataset
bytes, device, checkpoints, outputs, and completion status into a
verifiable receipt. A negative or null scientific result remains a valid
citable outcome when those implementation and provenance gates pass; it
blocks the intended positive claim,
\end{frame}

\begin{frame}{Future work: testing the\ldots{} (part 8)}
not the software release.
\end{frame}

\begin{frame}{Future work: testing the\ldots{} (part 9)}
A separate Friston protocol lane will resolve a machine-readable parity
matrix before reconstructing source experiments in Python. Until every
required source parameter, routine, unit, and estimand is recovered,
that lane will be described as a paper-constrained reconstruction rather
than an exact replication.
\end{frame}

\begin{frame}{Make the sharp server heuristic variational}
\protect\phantomsection\label{sec:future-server}
The asymmetry between the robustness axes
(sec.~3.3) motivates a consequential theoretical
design question that is explicitly parked outside the authorized
roadmap.
\end{frame}

\begin{frame}[fragile]{Make the sharp server\ldots{} (part 2)}
The client-side \(\beta\)/rcce update carries a derived, loss-specific
bounded-influence result under the matching assumptions; the sharp
server-side \texttt{robust\_aggregate} carries only its recovery limit
(eq.~10).
\end{frame}

\begin{frame}{Make the sharp server\ldots{} (part 3)}
The new variational aggregator (sec.~5.8.2,
sec.~12) supplies a \emph{rigorous} server-side
rule --- exact block updates descending the stated free energy
eq.~12 through non-increasing block updates,
\end{frame}

\begin{frame}{Make the sharp server\ldots{} (part 4)}
with a proven raw effective-weight bound and the same recovery corner.
What it costs is conservatism: it is the maximum-entropy consensus, not
the sharp accuracy-maximizer.
\end{frame}

\begin{frame}{Make the sharp server\ldots{} (part 5)}
That parked question is sharper than before: write down a generalized
variational objective in the FedGVI family (Mildner et al. 2025),
informed by recent closed-form GVI characterizations (Nguyen and
Westerhout 2026), logarithmic-pool weighting theory (Carvalho et al.
2023),
\end{frame}

\begin{frame}{Make the sharp server\ldots{} (part 6)}
and robust divergence-weighted federated aggregation (Li et al. 2022),
whose closed-form minimizer is competitive with the empirical
reweighting across declared contamination regimes.
\end{frame}

\begin{frame}{Make the sharp server\ldots{} (part 7)}
That would combine axis 3's effective-weight bound with axis 2's
empirical sharpness in one server rule. The recovery corner and the
variational objective together supply two boundary conditions any such
unification must satisfy.
\end{frame}

\begin{frame}[fragile]{Make the sharp server\ldots{} (part 8)}
Any empirical choice of \texttt{robustness} or \texttt{entropy\_weight}
will be made on separate calibration worlds using a proper log score,
then frozen before confirmatory evaluation. This guards against
selecting an apparent leader with evaluation truth and preserves null or
reversed confirmatory outcomes.
\end{frame}

\begin{frame}{Make the sharp server\ldots{} (part 9)}
The comparison family will include logarithmic and linear pools, the
current heuristic, the variational family, and a centered-log-ratio
geometric-median control; none inherits a parameter-space
robust-federated-learning guarantee merely by operating on beliefs.
\end{frame}

\begin{frame}{Promote the baseline to original FedGVI scale}
\protect\phantomsection\label{sec:future-scale}
The bounded-influence result is anchored locally by the small NumPy
logistic-regression baseline (fig.~18).
\end{frame}

\begin{frame}{Promote the baseline to\ldots{} (part 2)}
Promoting it to the GPU-scale Bayesian-neural-network experiments of the
source paper (Mildner et al. 2025) --- the experiments deferred here
(sec.~8.13) ---
\end{frame}

\begin{frame}{Promote the baseline to\ldots{} (part 3)}
would test whether the per-client robustness curve holds at the model
capacity and contamination regimes where federated learning (McMahan et
al. 2017) actually operates.
\end{frame}

\begin{frame}{Promote the baseline to\ldots{} (part 4)}
It would connect the discrete-POMDP result to the partitioned-VI line
(Ashman et al. 2022; Bui et al. 2018). The planned comparison would also
require posterior-parameterization parity with Bayesian neural-network
work, rather than treating the current deterministic point-mass MLP as a
posterior (Mildner et al. 2025).
\end{frame}

\begin{frame}{Promote the baseline to\ldots{} (part 5)}
The portable lane preserves the source protocol's site factors, client
cavity, factor-replacement update in natural coordinates. A synthetic
CPU/MPS pilot already exercises the cavity-conditioned local optimizer,
explicit device and fallback receipts, and checkpoint/resume
equivalence.
\end{frame}

\begin{frame}{Promote the baseline to\ldots{} (part 6)}
It does not establish source-dataset parity.
\end{frame}

\begin{frame}{Promote the baseline to\ldots{} (part 7)}
The next local campaign distinguishes a locked portable CPU/MPS profile
from an exact source-scale CUDA profile that remains external until
suitable hardware is available. FashionMNIST anchors source-dataset
protocol parity, while MNIST and KMNIST test portability.
\end{frame}

\begin{frame}{Promote the baseline to\ldots{} (part 8)}
A separate source-bound tabular pack will report proper-score effects
per licensed dataset, with training-only preprocessing and byte-,
split-, and license-level provenance; nested seeds will not be treated
as independent datasets.
\end{frame}

\begin{frame}{Extend hierarchical federation beyond the current stack}
\protect\phantomsection\label{sec:future-hierarchical}
The governing caveat here is that added depth must be shown to
\emph{earn} generalization, not merely to execute:
\end{frame}

\begin{frame}{Extend hierarchical\ldots{} (part 2)}
on the current sentinel task the deeper stacks match rather than beat
the flat baseline on location (sec.~8.15). The
extension therefore has two distinct fronts ---
\end{frame}

\begin{frame}{Extend hierarchical\ldots{} (part 3)}
carrying the recovery contract to deeper stacks, and finding a task
family in which depth actually pays.
\end{frame}

\begin{frame}[fragile]{Extend hierarchical\ldots{} (part 4)}
The 2-level hierarchical POMDP (sec.~14.9)
couples location inference to a single global context. The generic
N-level architecture uses \texttt{build\_nlevel\_world} from the
\texttt{fedference.pomdp} module and has already been exercised with a
3-level stack (sec.~14.11,
\end{frame}

\begin{frame}{Extend hierarchical\ldots{} (part 5)}
sec.~14.12): a meta-context variable (L3) gates the
context prior (L2) which in turn gates the location prior (L1).
\end{frame}

\begin{frame}{Extend hierarchical\ldots{} (part 6)}
The empirical-prior top-down messages (eq.~38,
eq.~39) remain valid variational steps at every
depth, and the log-linear-pool federation at each level is bit-identical
to the in-process result (Proposition
12).
\end{frame}

\begin{frame}{Extend hierarchical\ldots{} (part 7)}
The natural next question is whether the limit-as-proof contract of
sec.~8.13 survives still deeper hierarchies: does the
recovery corner (context prior \texorpdfstring{\ensuremath{\to}}{->} uniform) remain checkable to machine
precision when the L2 \texorpdfstring{\ensuremath{\to}}{->} L1 message is itself a function of an L3 belief
coupled to an L4 belief,
\end{frame}

\begin{frame}{Extend hierarchical\ldots{} (part 8)}
and can message-passing engines such as RxInfer (Bagaev et al. 2023)
carry the generic alternating-minimization at scale?
\end{frame}

\begin{frame}{Extend hierarchical\ldots{} (part 9)}
The current hierarchical diagnostic asks a narrower question. Under one
configured surprise threshold, it classifies the non-gating top level
below the threshold and the informative top level above it
(sec.~15.1). This is an implementation
control,
\end{frame}

\begin{frame}{Extend hierarchical\ldots{} (part 10)}
not a posterior refit or an autonomous depth-selection result.
\end{frame}

\begin{frame}{Extend hierarchical\ldots{} (part 11)}
A future breadth-and- depth search must define its candidate family,
scoring rule, uncertainty, replication units, and out-of-sample
evaluation before it can support a model- selection claim for the
generic N-level stack.
\end{frame}

\begin{frame}{Extend hierarchical\ldots{} (part 12)}
The next task family is deliberately controlled rather than merely
deeper: partially observable Four Rooms and Key-Door will compare flat,
oracle, learned, shuffled, and non-gating hierarchies at matched
horizons and compute. Task is the higher-level replication unit. A gain
in only one task will remain task-specific,
\end{frame}

\begin{frame}{Extend hierarchical\ldots{} (part 13)}
and the larger campaign will not begin until the hybrid representation
recovery gates pass.
\end{frame}

\begin{frame}{Move from process transport to true multi-machine
federation}
\protect\phantomsection\label{sec:future-transport}
The current deployment uses queue-backed single-machine processes and
loopback sockets. Promotion to cross-host workers would address the
remaining deployment caveat of sec.~8.13.
\end{frame}

\begin{frame}{Move from process\ldots{} (part 2)}
That extension must preserve the bit-identical consensus property proved
in Proposition 12.
\end{frame}

\begin{frame}[fragile]{Move from process\ldots{} (part 3)}
The \texttt{federation/} package and federation tests already establish
the API contract: a server collects serialized beliefs from 5 worker
channels, fuses them with \texttt{robust\_aggregate} at robustness
\(c = 1.5\), and broadcasts the consensus back over response channels,
with bit-identity verified at True.
\end{frame}

\begin{frame}{Move from process\ldots{} (part 4)}
The loopback socket path adds optional pre-shared-key frame integrity
and file-backed digest-verified replay validation. The next systems step
is an explicitly local Docker multi-node emulator with mTLS by default,
HMAC compatibility, checkpoint/restart, and reproducible message-fault
controls.
\end{frame}

\begin{frame}{Move from process\ldots{} (part 5)}
It is not physical multi-host evidence.

That later claim requires receipts from distinct hosts, deployment-grade
key management, timeout policy, and long-running orchestration across
process restarts, but it does not require changing the mathematics.
\end{frame}

\begin{frame}{Move from process\ldots{} (part 6)}
Secure aggregation, differential privacy, and Byzantine tolerance are
separate future threat models; transport integrity alone would not
establish any of them (Blanchard et al. 2017; Pillutla et al. 2022).
\end{frame}

\begin{frame}{Move beyond categorical state spaces}
\protect\phantomsection\label{sec:future-continuous}
This work is discrete-categorical, matching the community's worked POMDP
(Da Costa et al. 2020). A first step is already in place: the
closed-form Gaussian KL and Rényi divergences of
sec.~13 show the divergence family --- and its
\(\alpha\to1\) recovery --- carries over verbatim to Gaussian beliefs,
\end{frame}

\begin{frame}{Move beyond categorical\ldots{} (part 2)}
scoped out of the categorical experiments.
\end{frame}

\begin{frame}{Move beyond categorical\ldots{} (part 3)}
Continuous-state Gaussian generative models would test whether the
limit-as-proof contract survives the move off categorical state spaces
--- whether the recovery corner remains checkable to machine precision
when the belief simplex is replaced by a Gaussian belief,
\end{frame}

\begin{frame}{Move beyond categorical\ldots{} (part 4)}
and whether message-passing engines such as RxInfer (Bagaev et al. 2023)
can carry the robust update at scale.
\end{frame}

\begin{frame}{Move beyond categorical\ldots{} (part 5)}
Extending the structure-learning study (Smith et al. 2022; Friston and
Penny 2011) into continuous models would, in parallel, test whether
robust belief fusion and robust \emph{structure} fusion compose.
\end{frame}

\begin{frame}{Move beyond categorical\ldots{} (part 6)}
A minimal executable fixture now gates a discrete dynamics context over
continuous position and velocity, Gaussian observations, and bounded
actions. The bounded pilot now includes matched naive, robust,
discrete-only, continuous-only, and oracle-context components, a
singular-covariance rejection,
\end{frame}

\begin{frame}{Move beyond categorical\ldots{} (part 7)}
and next-position predictive scoring.
\end{frame}

\begin{frame}{Move beyond categorical\ldots{} (part 8)}
It is still a representation, recovery, and control surface rather than
confirmatory task evidence. The full study must freeze calibration and
budgets, use independently held-out worlds, retain outlier falsifiers,
and execute the preregistered comparison before supporting a hybrid-task
claim.
\end{frame}

\end{document}
